% --- Archon physics formalization source begin ---
% archon:physics
% archon:covers PhyXMiniProblems/problem_phyx_mini_0160.lean
% archon:source-report reports/phyx_mini/problem_phyx_mini_0160.source.json

\chapter{Physics problem phyx\_mini\_0160}
\label{ch:PhyXMiniProblems_problem_phyx_mini_0160}

\paragraph{Problem source.}
\#\# Physical scenario

An object is moved along the central axis of a spherical mirror while the lateral magnification \$m\$ of it is measured. Figure gives \$m\$ versus object distance \$p\$ for the range \$p\_a = 2.0 \textbackslash{}text\{cm\}\$ to \$p\_b = 8.0 \textbackslash{}text\{cm\}\$.

\#\# Question

What is \$m\$ for \$p = 14.0 \textbackslash{}text\{cm\}\$?

\#\# Answer choices

- A: A: -4

- B: B: -3.5

- C: C: -3

- D: D: -2.5

\#\# Figure caption (auxiliary; use the image as primary evidence)

The image is a graph with the horizontal axis labeled "p (cm)" and the vertical axis labeled "m". The x-axis ranges from \textbackslash{}( p\_a \textbackslash{}) to \textbackslash{}( p\_b \textbackslash{}), and the y-axis ranges from 0 to 4. 

There is a curved line on the graph that starts at a point near ( \textbackslash{}( p\_a \textbackslash{}), 2) and slopes upward, becoming steeper as it approaches ( \textbackslash{}( p\_b \textbackslash{}), 4). The line indicates a positive, nonlinear relationship between \textbackslash{}( p \textbackslash{}) and \textbackslash{}( m \textbackslash{}). Grid lines are present, helping to navigate the scale.

\#\# Dataset metadata

Geometrical Optics; Physical Model Grounding Reasoning, Multi-Formula Reasoning

\paragraph{Recorded answer/context.}
D: D: -2.5

\paragraph{Figure/image path.}
/root/proposal\_for\_physic/science-mango/phyx\_mini\_run/phyx\_data/test\_image/160.png

\paragraph{Formalization target.}
create a compiling Lean file with sorry bodies at `PhyXMiniProblems/problem\_phyx\_mini\_0160.lean`. The Lean declarations must preserve the physical quantities, dimensions or dimensional roles, figure labels, governing-law hypotheses, and final relation expressed by this problem.
Use Mathlib/Physlib names found through LeanExplore where available. If a physics API is missing, introduce faithful local abstractions rather than scalar placeholder aliases.

\begin{theorem}[Physics formalization target]
\label{thm:physics:phyx_mini_0160:target}
\lean{PhyXMiniProblems.ProblemPhyXMini0160.problem_phyx_mini_0160}
\uses{def:physics:phyx-mini-0160:phyxminiproblems-problemphyxmini0160-sphericalmirrormagnificationsetup, def:physics:phyx-mini-0160:phyxminiproblems-problemphyxmini0160-matchesproblemandgraph, def:physics:phyx-mini-0160:phyxminiproblems-problemphyxmini0160-hasfigurederivedfocalcalibration, def:physics:phyx-mini-0160:phyxminiproblems-problemphyxmini0160-hasphysicalparameters, def:physics:phyx-mini-0160:phyxminiproblems-problemphyxmini0160-satisfiesparaxialsphericalmirrorlaws, def:physics:phyx-mini-0160:phyxminiproblems-problemphyxmini0160-recordeddatasetanswer, def:physics:phyx-mini-0160:phyxminiproblems-problemphyxmini0160-matchesdisplayedanswer}
The assigned autoformalize agent should translate this physics problem into Lean declarations in the covered file, with theorem and lemma proof bodies written as `by sorry`.
\end{theorem}
\begin{proof}
This is an autoformalization task, not a proof task. Produce statements that can later be proved by the physics prover without weakening the physical model.
\end{proof}

% --- Archon declaration topology begin ---
\section{Declaration topology}
\begin{definition}[Length Quantity]
  \label{def:physics:phyx-mini-0160:phyxminiproblems-problemphyxmini0160-lengthquantity}
  \lean{PhyXMiniProblems.ProblemPhyXMini0160.LengthQuantity}
  A signed physical length represented coherently in every choice of units
\end{definition}
\begin{proof}
  This is the declared model definition of Length Quantity.
\end{proof}
\begin{definition}[length Readout]
  \label{def:physics:phyx-mini-0160:phyxminiproblems-problemphyxmini0160-lengthreadout}
  \lean{PhyXMiniProblems.ProblemPhyXMini0160.lengthReadout}
  \uses{def:physics:phyx-mini-0160:phyxminiproblems-problemphyxmini0160-lengthquantity}
  Scalar readout of a signed physical length in a selected length unit
\end{definition}
\begin{proof}
  This is the declared model definition of length Readout.
\end{proof}
\begin{definition}[length In Centimeters]
  \label{def:physics:phyx-mini-0160:phyxminiproblems-problemphyxmini0160-lengthincentimeters}
  \lean{PhyXMiniProblems.ProblemPhyXMini0160.lengthInCentimeters}
  \uses{def:physics:phyx-mini-0160:phyxminiproblems-problemphyxmini0160-lengthquantity, def:physics:phyx-mini-0160:phyxminiproblems-problemphyxmini0160-lengthreadout}
  The centimeter readout used by the graph and the question
\end{definition}
\begin{proof}
  This is the declared model definition of length In Centimeters.
\end{proof}
\begin{definition}[Spherical Mirror Kind]
  \label{def:physics:phyx-mini-0160:phyxminiproblems-problemphyxmini0160-sphericalmirrorkind}
  \lean{PhyXMiniProblems.ProblemPhyXMini0160.SphericalMirrorKind}
  The two spherical-mirror orientations distinguished by their focal signs
\end{definition}
\begin{proof}
  This is the declared model definition of Spherical Mirror Kind.
\end{proof}
\begin{definition}[Graph Variable]
  \label{def:physics:phyx-mini-0160:phyxminiproblems-problemphyxmini0160-graphvariable}
  \lean{PhyXMiniProblems.ProblemPhyXMini0160.GraphVariable}
  Physical quantities plotted on the two axes of the supplied graph
\end{definition}
\begin{proof}
  This is the declared model definition of Graph Variable.
\end{proof}
\begin{definition}[Spherical Mirror Magnification Setup]
  \label{def:physics:phyx-mini-0160:phyxminiproblems-problemphyxmini0160-sphericalmirrormagnificationsetup}
  \lean{PhyXMiniProblems.ProblemPhyXMini0160.SphericalMirrorMagnificationSetup}
  \uses{def:physics:phyx-mini-0160:phyxminiproblems-problemphyxmini0160-lengthquantity, def:physics:phyx-mini-0160:phyxminiproblems-problemphyxmini0160-sphericalmirrorkind, def:physics:phyx-mini-0160:phyxminiproblems-problemphyxmini0160-graphvariable}
  The mirror, its signed paraxial image-distance response, and the measured magnification curve. The three named object distances are the graph endpoints p\_a, p\_b, and the requested distance p = 14 cm. No field fixes the magnification at the requested distance
\end{definition}
\begin{proof}
  This is the declared model definition of Spherical Mirror Magnification Setup.
\end{proof}
\begin{definition}[Matches Problem And Graph]
  \label{def:physics:phyx-mini-0160:phyxminiproblems-problemphyxmini0160-matchesproblemandgraph}
  \lean{PhyXMiniProblems.ProblemPhyXMini0160.MatchesProblemAndGraph}
  \uses{def:physics:phyx-mini-0160:phyxminiproblems-problemphyxmini0160-lengthquantity, def:physics:phyx-mini-0160:phyxminiproblems-problemphyxmini0160-lengthincentimeters, def:physics:phyx-mini-0160:phyxminiproblems-problemphyxmini0160-sphericalmirrormagnificationsetup}
  The explicit problem and intended-graph readouts. Besides the axis labels and the object-distance range, the source caption records a positive, increasing branch of the magnification curve on [p\_a,p\_b]
\end{definition}
\begin{proof}
  This is the declared model definition of Matches Problem And Graph.
\end{proof}
\begin{definition}[Has Figure Derived Focal Calibration]
  \label{def:physics:phyx-mini-0160:phyxminiproblems-problemphyxmini0160-hasfigurederivedfocalcalibration}
  \lean{PhyXMiniProblems.ProblemPhyXMini0160.HasFigureDerivedFocalCalibration}
  \uses{def:physics:phyx-mini-0160:phyxminiproblems-problemphyxmini0160-lengthincentimeters, def:physics:phyx-mini-0160:phyxminiproblems-problemphyxmini0160-sphericalmirrormagnificationsetup}
  The intended magnification curve's spherical-mirror calibration is a focal length of 10 cm. It is kept as a separate intermediate graph/model input, not as the requested magnification at 14 cm
\end{definition}
\begin{proof}
  This is the declared model definition of Has Figure Derived Focal Calibration.
\end{proof}
\begin{definition}[Has Physical Parameters]
  \label{def:physics:phyx-mini-0160:phyxminiproblems-problemphyxmini0160-hasphysicalparameters}
  \lean{PhyXMiniProblems.ProblemPhyXMini0160.HasPhysicalParameters}
  \uses{def:physics:phyx-mini-0160:phyxminiproblems-problemphyxmini0160-lengthincentimeters, def:physics:phyx-mini-0160:phyxminiproblems-problemphyxmini0160-sphericalmirrormagnificationsetup}
  Positivity and finite-distance conditions for the stated configuration
\end{definition}
\begin{proof}
  This is the declared model definition of Has Physical Parameters.
\end{proof}
\begin{definition}[Is Finite Object Configuration]
  \label{def:physics:phyx-mini-0160:phyxminiproblems-problemphyxmini0160-isfiniteobjectconfiguration}
  \lean{PhyXMiniProblems.ProblemPhyXMini0160.IsFiniteObjectConfiguration}
  \uses{def:physics:phyx-mini-0160:phyxminiproblems-problemphyxmini0160-lengthquantity, def:physics:phyx-mini-0160:phyxminiproblems-problemphyxmini0160-lengthincentimeters, def:physics:phyx-mini-0160:phyxminiproblems-problemphyxmini0160-sphericalmirrormagnificationsetup}
  A positive object distance at which the paraxial image is finite
\end{definition}
\begin{proof}
  This is the declared model definition of Is Finite Object Configuration.
\end{proof}
\begin{definition}[Satisfies Paraxial Spherical Mirror Laws]
  \label{def:physics:phyx-mini-0160:phyxminiproblems-problemphyxmini0160-satisfiesparaxialsphericalmirrorlaws}
  \lean{PhyXMiniProblems.ProblemPhyXMini0160.SatisfiesParaxialSphericalMirrorLaws}
  \uses{def:physics:phyx-mini-0160:phyxminiproblems-problemphyxmini0160-lengthquantity, def:physics:phyx-mini-0160:phyxminiproblems-problemphyxmini0160-lengthreadout, def:physics:phyx-mini-0160:phyxminiproblems-problemphyxmini0160-sphericalmirrormagnificationsetup, def:physics:phyx-mini-0160:phyxminiproblems-problemphyxmini0160-isfiniteobjectconfiguration}
  The paraxial spherical-mirror equation 1/p + 1/q = 1/f and the signed lateral-magnification law m = -q/p. Both laws are stated in every length unit, so the physical content is not tied to the centimeter readouts used by the graph
\end{definition}
\begin{proof}
  This is the declared model definition of Satisfies Paraxial Spherical Mirror Laws.
\end{proof}
\begin{definition}[Answer Choice]
  \label{def:physics:phyx-mini-0160:phyxminiproblems-problemphyxmini0160-answerchoice}
  \lean{PhyXMiniProblems.ProblemPhyXMini0160.AnswerChoice}
  Labels of the four answer choices printed in the source problem
\end{definition}
\begin{proof}
  This is the declared model definition of Answer Choice.
\end{proof}
\begin{definition}[displayed Magnification]
  \label{def:physics:phyx-mini-0160:phyxminiproblems-problemphyxmini0160-answerchoice-displayedmagnification}
  \lean{PhyXMiniProblems.ProblemPhyXMini0160.AnswerChoice.displayedMagnification}
  \uses{def:physics:phyx-mini-0160:phyxminiproblems-problemphyxmini0160-answerchoice}
  Dimensionless magnification displayed beside each answer choice
\end{definition}
\begin{proof}
  This is the declared model definition of displayed Magnification.
\end{proof}
\begin{definition}[recorded Dataset Answer]
  \label{def:physics:phyx-mini-0160:phyxminiproblems-problemphyxmini0160-recordeddatasetanswer}
  \lean{PhyXMiniProblems.ProblemPhyXMini0160.recordedDatasetAnswer}
  \uses{def:physics:phyx-mini-0160:phyxminiproblems-problemphyxmini0160-answerchoice}
  The answer label recorded in the supplied dataset
\end{definition}
\begin{proof}
  This is the declared model definition of recorded Dataset Answer.
\end{proof}
\begin{definition}[Matches Displayed Answer]
  \label{def:physics:phyx-mini-0160:phyxminiproblems-problemphyxmini0160-matchesdisplayedanswer}
  \lean{PhyXMiniProblems.ProblemPhyXMini0160.MatchesDisplayedAnswer}
  \uses{def:physics:phyx-mini-0160:phyxminiproblems-problemphyxmini0160-sphericalmirrormagnificationsetup, def:physics:phyx-mini-0160:phyxminiproblems-problemphyxmini0160-answerchoice}
  A choice agrees exactly with the requested dimensionless magnification
\end{definition}
\begin{proof}
  This is the declared model definition of Matches Displayed Answer.
\end{proof}
\begin{lemma}[requested Signed Image Distance Centimeters eq thirty five]
  \label{lem:physics:phyx-mini-0160:phyxminiproblems-problemphyxmini0160-requestedsignedimagedistancecentimeters-eq-thirty-five}
  \lean{PhyXMiniProblems.ProblemPhyXMini0160.requestedSignedImageDistanceCentimeters_eq_thirty_five}
  \uses{def:physics:phyx-mini-0160:phyxminiproblems-problemphyxmini0160-lengthincentimeters, def:physics:phyx-mini-0160:phyxminiproblems-problemphyxmini0160-sphericalmirrormagnificationsetup, def:physics:phyx-mini-0160:phyxminiproblems-problemphyxmini0160-matchesproblemandgraph, def:physics:phyx-mini-0160:phyxminiproblems-problemphyxmini0160-hasfigurederivedfocalcalibration, def:physics:phyx-mini-0160:phyxminiproblems-problemphyxmini0160-hasphysicalparameters, def:physics:phyx-mini-0160:phyxminiproblems-problemphyxmini0160-satisfiesparaxialsphericalmirrorlaws}
  At p = 14 cm and f = 10 cm, the mirror equation gives the intermediate signed image distance q = 35 cm
\end{lemma}
\begin{proof}
  The conclusion follows from the governing relations and figure readouts named in the statement of requested Signed Image Distance Centimeters eq thirty five.
\end{proof}
% --- Archon declaration topology end ---
% --- Archon physics formalization source end ---
